\section{Discussion}
\label{sec:discussion}

\paragraph{Mask-based architectures help anomaly detection.}
\todo{Summarize the central empirical finding.}  
% e.g., "Our results show that EoMT consistently outperforms ERFNet on most 
% anomaly benchmarks, supporting the hypothesis that mask classification 
% provides a stronger inductive bias for detecting out-of-distribution regions."

\todo{Explain why, briefly.}  
% e.g., "Because each query in EoMT learns a near-independent one-vs-all 
% classifier, regions claimed by no query produce a clean rejection signal — 
% something a per-pixel softmax model cannot replicate."

\paragraph{Fine-tuning recovers most of the closed-set gap.}
\todo{State the main fine-tuning finding.}  
% "Fine-tuning EoMT-COCO on Cityscapes closes most of the mIoU gap to native 
% training (XX.XX\% vs YY.YY\%) and brings anomaly detection performance to 
% a comparable level, even with a frozen backbone and 20 epochs of training."

\todo{Comment on efficiency.}  
% "This suggests that broad pre-training on COCO transfers efficiently to 
% road scenes, and full backbone updates are not required."

\paragraph{Limitations.}
\todo{Acknowledge limits.}  
% Two or three things: small compute budget (Colab L4, ~20 epochs only); 
% smaller evaluation subsets than the full benchmarks; no outlier-supervised 
% baselines (PEBAL, DenseHybrid) for comparison. RbA has high FPR@95 trade-off.

\paragraph{Future work.}
\todo{One or two directions.}  
% e.g., "Extending fine-tuning to unfrozen backbone training, evaluating 
% on the full benchmark splits, and combining post-hoc scores with light 
% outlier exposure are natural next steps."